% GENERATED FILE. Edit canonical lessons, then run npm run generate:books.
% canonical-source-hash: fnv1a64:34c397d3fc933a66
% canonical-lessons: MW-C138-balad, MW-C138-bachhro, MW-C138-hiran, MW-C138-lukri, MW-C138-siyal

\chapter{Ox, Calf, Deer, Fox, Jackal}
\label{ch:mw-ox-calf-deer-fox-jackal}
\hlchaptermodality{138}

\begin{chapteropening}
\textbf{By the end of this chapter:} \emph{I can say an ox, a calf, a deer, a fox and a jackal.}

Complete the last lesson of chapter 138: I can say an ox, a calf, a deer, a fox and a jackal.
\end{chapteropening}

\section[\texorpdfstring{\mw{बळद} (baḷad)}{बळद (baḷad)}]{\mw{बळद} (baḷad) — an ox}

\label{lesson:MW-C138-balad}



\begin{quote}
Before the new one: say the Marwadi for dates, then the Marwadi for a coconut.
\end{quote}

\subsection*{You'll want to know: \mw{बळद}}
\textbf{\mw{बळद}} — \emph{baḷad} — \textquotedblleft{}an ox\textquotedblright{}.

\begin{cousinweb}
\textbf{\mw{बळद}} (\emph{baḷad}) is a tadbhava, a word inherited from Sanskrit and reshaped by sound change. It comes from \emph{balivarda}, \textquotedblleft{}ox, bull\textquotedblright{}. CDIAL gives no Marwari form in this entry (CDIAL 9176); the match is through Gujarati \emph{baḷad} and Hindi \emph{balad}. CDIAL records a view that the Sanskrit word joins \emph{balin}, \textquotedblleft{}strong\textquotedblright{}, with a non-Aryan word for \textquotedblleft{}ox\textquotedblright{}.
\end{cousinweb}

\begin{grammarlens}[title={What you've built}]
One more word for this chapter.
\end{grammarlens}

\subsection*{Your turn}
\begin{itemize}
\raggedright
  \item \emph{Say it:} \emph{baḷad}
  \item \emph{Say it:} \emph{baḷad}, once more
  \item \emph{Say it:} say \emph{baḷad}
  \item \emph{Recall:} say the Marwadi for dates, then the Marwadi for a coconut, then say \emph{baḷad} again
\end{itemize}

\begin{tcolorbox}[breakable,colback=teal!4,colframe=teal!35!black,title={Before you move on}]
What does \mw{बळद} mean? (\textquotedblleft{}An ox\textquotedblright{}.) Say it once more.
\end{tcolorbox}
\section[\texorpdfstring{\mw{बाछड़ो} (bāchhṛo)}{बाछड़ो (bāchhṛo)}]{\mw{बाछड़ो} (bāchhṛo) — a calf}

\label{lesson:MW-C138-bachhro}



\begin{quote}
Before the new one: say the Marwadi for a coconut, then the Marwadi for an ox.
\end{quote}

\subsection*{You'll want to know: \mw{बाछड़ो}}
\textbf{\mw{बाछड़ो}} — \emph{bāchhṛo} — \textquotedblleft{}a calf\textquotedblright{}.

\begin{cousinweb}
\textbf{\mw{बाछड़ो}} (\emph{bāchhṛo}) is a tadbhava, a word inherited from Sanskrit and reshaped by sound change. It comes from \emph{vatsa}, \textquotedblleft{}calf, child\textquotedblright{}. CDIAL gives no Marwari form in this entry (CDIAL 11239); the match is through Hindi \emph{bāchā}. The \emph{ṛ} is an added suffix. The Sanskrit word first meant \textquotedblleft{}yearling\textquotedblright{}.
\end{cousinweb}

\begin{grammarlens}[title={What you've built}]
One more word for this chapter.
\end{grammarlens}

\subsection*{Your turn}
\begin{itemize}
\raggedright
  \item \emph{Say it:} \emph{bāchhṛo}
  \item \emph{Say it:} \emph{bāchhṛo}, once more
  \item \emph{Say it:} say \emph{bāchhṛo}
  \item \emph{Recall:} say the Marwadi for a coconut, then the Marwadi for an ox, then say \emph{bāchhṛo} again
\end{itemize}

\begin{tcolorbox}[breakable,colback=teal!4,colframe=teal!35!black,title={Before you move on}]
What does \mw{बाछड़ो} mean? (\textquotedblleft{}A calf\textquotedblright{}.) Say it once more.
\end{tcolorbox}
\section[\texorpdfstring{\mw{हिरण} (hiraṇ)}{हिरण (hiraṇ)}]{\mw{हिरण} (hiraṇ) — a deer}

\label{lesson:MW-C138-hiran}



\begin{quote}
Before the new one: say the Marwadi for an ox, then the Marwadi for a calf.
\end{quote}

\subsection*{You'll want to know: \mw{हिरण}}
\textbf{\mw{हिरण}} — \emph{hiraṇ} — \textquotedblleft{}a deer\textquotedblright{}.

\begin{cousinweb}
\textbf{\mw{हिरण}} (\emph{hiraṇ}) is a tadbhava, a word inherited from Sanskrit and reshaped by sound change. It comes from \emph{hariṇa}, \textquotedblleft{}deer\textquotedblright{}. CDIAL gives no Marwari form in this entry (CDIAL 13982); the match is through Hindi \emph{hiran} and Gujarati \emph{haraṇ}. The Sanskrit word first meant \textquotedblleft{}yellow\textquotedblright{}. CDIAL rejects a proposed Austroasiatic source.
\end{cousinweb}

\begin{grammarlens}[title={What you've built}]
One more word for this chapter.
\end{grammarlens}

\subsection*{Your turn}
\begin{itemize}
\raggedright
  \item \emph{Say it:} \emph{hiraṇ}
  \item \emph{Say it:} \emph{hiraṇ}, once more
  \item \emph{Say it:} say \emph{hiraṇ}
  \item \emph{Recall:} say the Marwadi for an ox, then the Marwadi for a calf, then say \emph{hiraṇ} again
\end{itemize}

\begin{tcolorbox}[breakable,colback=teal!4,colframe=teal!35!black,title={Before you move on}]
What does \mw{हिरण} mean? (\textquotedblleft{}A deer\textquotedblright{}.) Say it once more.
\end{tcolorbox}
\section[\texorpdfstring{\mw{लूंकड़ी} (lū̃kṛī)}{लूंकड़ी (lū̃kṛī)}]{\mw{लूंकड़ी} (lū̃kṛī) — a fox}

\label{lesson:MW-C138-lukri}



\begin{quote}
Before the new one: say the Marwadi for a calf, then the Marwadi for a deer.
\end{quote}

\subsection*{You'll want to know: \mw{लूंकड़ी}}
\textbf{\mw{लूंकड़ी}} — \emph{lū̃kṛī} — \textquotedblleft{}a fox\textquotedblright{}.

\begin{cousinweb}
The origin of \textbf{\mw{लूंकड़ी}} (\emph{lū̃kṛī}) is not settled. CDIAL lists Hindi \emph{lomṛī}, \textquotedblleft{}fox\textquotedblright{}, under Sanskrit \emph{lomaṭaka}, \textquotedblleft{}fox\textquotedblright{} (CDIAL 11153), but no form with the \emph{lū̃k-} of \textbf{\mw{लूंकड़ी}} (\emph{lū̃kṛī}).
\end{cousinweb}

\begin{grammarlens}[title={What you've built}]
One more word for this chapter.
\end{grammarlens}

\subsection*{Your turn}
\begin{itemize}
\raggedright
  \item \emph{Say it:} \emph{lū̃kṛī}
  \item \emph{Say it:} \emph{lū̃kṛī}, once more
  \item \emph{Say it:} say \emph{lū̃kṛī}
  \item \emph{Recall:} say the Marwadi for a calf, then the Marwadi for a deer, then say \emph{lū̃kṛī} again
\end{itemize}

\begin{tcolorbox}[breakable,colback=teal!4,colframe=teal!35!black,title={Before you move on}]
What does \mw{लूंकड़ी} mean? (\textquotedblleft{}A fox\textquotedblright{}.) Say it once more.
\end{tcolorbox}
\section[\texorpdfstring{\mw{सियाळ} (siyāḷ)}{सियाळ (siyāḷ)}]{\mw{सियाळ} (siyāḷ) — a jackal}

\label{lesson:MW-C138-siyal}



\begin{quote}
Before the new one: say the Marwadi for a deer, then the Marwadi for a fox.
\end{quote}

\subsection*{You'll want to know: \mw{सियाळ}}
\textbf{\mw{सियाळ}} — \emph{siyāḷ} — \textquotedblleft{}a jackal\textquotedblright{}.

\begin{cousinweb}
\textbf{\mw{सियाळ}} (\emph{siyāḷ}) is a tadbhava, a word inherited from Sanskrit and reshaped by sound change. It comes from \emph{śṛgāla}, \textquotedblleft{}jackal\textquotedblright{}. CDIAL gives no Marwari form in this entry (CDIAL 12578); the match is through Hindi \emph{siyāl} and Gujarati \emph{siyāḷ}.
\end{cousinweb}

\begin{grammarlens}[title={What you've built}]
One more word for this chapter.
\end{grammarlens}

\subsection*{Your turn}
\begin{itemize}
\raggedright
  \item \emph{Say it:} \emph{siyāḷ}
  \item \emph{Say it:} \emph{siyāḷ}, once more
  \item \emph{Say it:} say \emph{siyāḷ}
  \item \emph{Recall:} say the Marwadi for a deer, then the Marwadi for a fox, then say \emph{siyāḷ} again
\end{itemize}

\begin{tcolorbox}[breakable,colback=teal!4,colframe=teal!35!black,title={Before you move on}]
What does \mw{सियाळ} mean? (\textquotedblleft{}A jackal\textquotedblright{}.) Say it once more.
\end{tcolorbox}
